% question.tex — Fragmento por questão (recriação visual do loop de questões
% de generatePdfWithDoc, supabase/functions/generate-exam-pdf/legacyPdfLib.ts:525-578).
%
% Referência de estrutura/estilo, não um arquivo literalmente \input'ado sem
% modificação: cada questão tem de 0 a 4 opções e uma imagem opcional, então o
% TS (Task 9, renderTemplate.ts) monta este fragmento dinamicamente por questão
% a partir de um template mais granular do que o pseudo-código abaixo sugere.
% O mecanismo exato de composição (substituição literal deste arquivo vs. TS
% montando a string seguindo este padrão) é decisão da Task 9, não desta task.
%
% Placeholders literais @@...@@, substituídos pelo TS (Task 9), repetido N
% vezes — um bloco destes por questão, sem loop LaTeX.

\needspace{4\baselineskip}
\vspace{12pt}
\textbf{\textcolor{winemid}{Questão @@Q_NUMBER@@}}\\[4pt]
@@Q_TEXT@@

@@Q_IMAGE_BLOCK@@

\begin{enumerate}[label=\textbf{@@Q_OPT_LABEL@@)}, leftmargin=*, itemsep=4pt]
@@Q_OPTIONS@@
\end{enumerate}

% ── Bloco de imagem (usado no lugar de @@Q_IMAGE_BLOCK@@ quando a questão tem
% imagem; substituído por string vazia quando não tem). Replica a heurística
% de tamanho de legacyPdfLib.ts:552-565 (teto de 82% da largura de conteúdo e
% 320pt de altura, o menor dos dois limites vence preservando proporção) —
% em LaTeX isso é só `width=...,height=...,keepaspectratio`, sem precisar
% calcular manualmente o `min(...)` como no pdf-lib original. graphicx também
% cuida da quebra de página da imagem em si automaticamente; o `\needspace`
% no topo deste arquivo cobre apenas o início do bloco da questão (equivalente
% ao `ensureSpace()` do motor antigo, legacyPdfLib.ts:521-523), garantindo que
% ao menos o título "Questão N" e a primeira linha do enunciado não fiquem
% órfãos no rodapé de uma página.
%
% \begin{center}
% \includegraphics[width=0.82\textwidth,height=320pt,keepaspectratio]{@@Q_IMAGE_PATH@@}
% \end{center}

% ── Padrão de uma opção dentro de @@Q_OPTIONS@@ (uma linha \item por opção,
% rótulo A-D já vem em maiúsculo do banco — ver OptionRow.label em types.ts) ──
% \item @@Q_OPT_TEXT@@
